\documentclass[11pt]{article}

\usepackage[margin=1in]{geometry}
\usepackage{amsmath,amssymb,amsthm,mathtools}
\usepackage{microtype}
\usepackage[hidelinks]{hyperref}
\usepackage{enumitem}

\newtheorem{theorem}{Theorem}
\newtheorem{lemma}[theorem]{Lemma}
\newtheorem{corollary}[theorem]{Corollary}
\newtheorem{proposition}[theorem]{Proposition}
\theoremstyle{remark}
\newtheorem{remark}[theorem]{Remark}
\newtheorem{example}[theorem]{Example}

\newcommand{\ind}[1]{\mathbf{1}\!\left\{#1\right\}}
\newcommand{\E}{\mathbb{E}}
\newcommand{\R}{\mathbb{R}}
\newcommand{\N}{\mathbb{N}}
\newcommand{\G}{\mathcal{G}}
\newcommand{\C}{\mathcal{C}}
\newcommand{\Reg}{\mathfrak{R}}
\newcommand{\Tau}{\tau}
\newcommand{\GammaTail}{\Gamma}

\title{Delays adversarial}
\author{}
\date{}

\begin{document}
\maketitle

\section{Model}

Fix a horizon $T\in\N$.  Before the game begins, an oblivious adversary chooses a
deterministic sequence
\[
    p=(p_1,\ldots,p_T)\in[0,1]^T.
\]
At every round $t=1,\ldots,T$:
\begin{enumerate}[label=(\arabic*)]
    \item the learner observes $p_t$;
    \item the learner posts a possibly randomized quote $A_t\in[0,1]$ based only
    on $p_1,\ldots,p_t$ and its internal randomization.
\end{enumerate}
The learner's actions do not affect the price sequence.

For a deterministic quote $a\in[0,1]$ posted at time~$t$, define its execution
delay by
\[
    \Tau_t^p(a)
    :=
    \min\bigl\{d\in\{1,\ldots,T-t\}:p_{t+d}>a\bigr\},
\]
with the convention $\min\varnothing=\infty$.  Its reward is defined as
\[
    r_t^p(a)
    :=
    a\,\ind{\Tau_t^p(a)<\infty}.
\]
Thus, the order executes at the first strictly larger future price and earns its
posted quote.

A randomized policy $\pi$ is a collection of nonanticipating decision rules for
$A_1,\ldots,A_T$.  Its regret on $p$ is
\begin{equation}\label{eq:regret}
    R_T(\pi,p)
    :=
    \sup_{a\in[0,1]}
    \sum_{t=1}^T r_t^p(a)
    -
    \E_\pi\!\left[\sum_{t=1}^T r_t^p(A_t)\right],
\end{equation}
where the expectation is only over the learner's randomization.

\section{Reward-weighted delay tails}

For an integer truncation level $H\ge 0$, define the truncated reward
\[
    r_{t,H}^p(a)
    :=
    a\,\ind{\Tau_t^p(a)\le H}.
\]
The reward lost by truncating quote $a$ after $H$ rounds is therefore
\[
    \sum_{t=1}^T\bigl(r_t^p(a)-r_{t,H}^p(a)\bigr)
    =
    \sum_{t=1}^T
    a\,\ind{H<\Tau_t^p(a)<\infty}.
\]
This motivates the cumulative reward-weighted tail
\begin{equation}\label{eq:Gamma}
    \GammaTail_{T,H}(p)
    :=
    \sup_{a\in[0,1]}
    \sum_{t=1}^T
    a\,\ind{H<\Tau_t^p(a)<\infty}.
\end{equation}
It measures the largest total reward that a single fixed quote obtains from
orders whose execution delay is strictly larger than~$H$.

Let
\[
    g:\N_0\longrightarrow[0,1]
\]
be nonincreasing, where $\N_0=\{0,1,2,\ldots\}$.  For each horizon~$T$, define
the maximal tail-envelope class of price sequences
\begin{equation}\label{eq:class}
    \C_T(g)
    :=
    \left\{
        p\in[0,1]^T:
        \GammaTail_{T,H}(p)\le T g(H)
        \text{ for every }H\in\N_0
    \right\}.
\end{equation}
We restrict attention to adversaries that select a sequence from this class.
The learner knows $T$ and $g$, but not the sequence selected from $\C_T(g)$.
Define the minimax regret
\begin{equation}\label{eq:minimax}
    \Reg_T(g)
    :=
    \inf_{\pi}
    \sup_{p\in\C_T(g)}R_T(\pi,p).
\end{equation}

\subsection*{Examples of tail envelopes}

Every price sequence satisfies $\GammaTail_{T,H}(p)\le T$, and hence belongs
to $\C_T(g)$ for the uninformative envelope $g\equiv 1$.  The following
examples give more informative envelopes.

\paragraph{Uniformly fast execution.}
Suppose that, for some $D\ge 1$, every order that eventually executes does so
within $D$ rounds:
\[
    \Tau_t^p(a)<\infty
    \quad\Longrightarrow\quad
    \Tau_t^p(a)\le D
    \qquad\text{for every $t$ and $a$.}
\]
Then $\GammaTail_{T,H}(p)=0$ for $H\ge D$, so $p\in\C_T(g_D)$ for
$g_D(H):=\ind{H<D}$.  In particular, every nonincreasing price sequence has
this property with $D=1$: if some later price exceeds $a$, then already
$p_{t+1}>a$.

\paragraph{Recurring spikes.}
Let $p_t\in\{0,1\}$, assume $p_1=1$, and write
$1=s_1<s_2<\cdots<s_m$ for the times at which $p_t=1$.  If
$d_i:=s_{i+1}-s_i$, then
\[
    \GammaTail_{T,H}(p)
    =
    \sum_{i=1}^{m-1}(d_i-H)_+,
    \qquad
    (u)_+:=\max\{u,0\}.
\]
Indeed, within a gap of length $d_i$, exactly $(d_i-H)_+$ orders wait more
than $H$ rounds for the next spike; the supremum in \eqref{eq:Gamma} is
approached as $a\uparrow 1$.  Thus regularly spaced spikes with gap $D$
satisfy
\[
    \GammaTail_{T,H}(p)
    \le
    T\left(1-\frac{H}{D}\right)_+,
\]
so one may take $g(H)=(1-H/D)_+$.  More generally, long gaps are allowed
provided that they occupy a suitably small fraction of time.

\paragraph{A single terminal jump.}
For $p=(0,\ldots,0,x)$ with $x\in(0,1]$,
\[
    \GammaTail_{T,H}(p)
    =
    x(T-H-1)_+.
\]
Consequently, this sequence belongs to $\C_T(g_x)$ for the constant envelope
$g_x(H)\equiv x$.  Although any one finite path has zero tail once
$H\ge T-1$, the family of terminal-jump paths across all horizons cannot be
covered by a single envelope that tends to zero.

\section{Main characterization}

\begin{theorem}[Adversarial delay-tail characterization]\label{thm:main}
Let $g:\N_0\to[0,1]$ be nonincreasing, and let $\Reg_T(g)$ be defined by
\eqref{eq:minimax}.

\begin{enumerate}[label=(\roman*)]
    \item For every $T\in\N$, every integer $H\ge0$, and every integer $K\ge2$,
    \begin{equation}\label{eq:upper-main}
        \Reg_T(g)
        \le
        \sqrt{\frac{T(H+1)\log K}{2}}
        +\frac{T}{K}
        +Tg(H).
    \end{equation}

    \item Let
    \[
        c:=\lim_{H\to\infty}g(H),
    \]
    which exists because $g$ is nonincreasing.  If $c>0$, then for every
    $T\in\N$,
    \begin{equation}\label{eq:lower-main}
        \Reg_T(g)
        \ge
        \frac{c}{e}(T-1).
    \end{equation}
\end{enumerate}
Consequently,
\begin{equation}\label{eq:iff}
    \boxed{
    \Reg_T(g)=o(T)
    \quad\Longleftrightarrow\quad
    \lim_{H\to\infty}g(H)=0.
    }
\end{equation}
\end{theorem}

\subsection*{Why the theorem is true}

If $g(H)$ is small, then no fixed quote can obtain much total reward from orders
that remain unresolved for more than~$H$ rounds.  We may therefore truncate all
orders after $H$ rounds at a cost of at most $Tg(H)$ against the best fixed
comparator.  After truncation, the complete reward vector for round~$t$ is known
by time $t+H$.  Splitting time into $H+1$ residue classes eliminates this delay:
before the next decision in a residue class, the preceding reward vector in that
class is already known.  Hedge can then be run separately on each class.  Finally,
rounding every comparator quote downward to a grid loses at most $1/K$ per round.

For the converse, suppose $g(H)$ remains bounded below by $c>0$.  The class then
contains every terminal-jump sequence
\[
    (0,\ldots,0,x),\qquad 0\le x\le c.
\]
Before the final round, all such sequences have exactly the same observed prefix,
so the learner must post its first $T-1$ quotes without knowing~$x$.  A suitable
distribution over terminal values creates an unavoidable loss of $c/e$ per
preterminal round relative to the best fixed quote in hindsight.

The proof is divided into an upper-bound part and a lower-bound part.

\section{Upper bound}

\subsection{Truncation}

\begin{lemma}[Pathwise truncation comparison]\label{lem:truncation}
Fix $T,H$, a sequence $p\in[0,1]^T$, and a policy $\pi$.  Define the truncated
regret
\[
    R_{T,H}(\pi,p)
    :=
    \sup_{a\in[0,1]}
    \sum_{t=1}^T r_{t,H}^p(a)
    -
    \E_\pi\!\left[\sum_{t=1}^T r_{t,H}^p(A_t)\right].
\]
Then
\begin{equation}\label{eq:truncation}
    R_T(\pi,p)
    \le
    R_{T,H}(\pi,p)+\GammaTail_{T,H}(p).
\end{equation}
\end{lemma}

\begin{proof}
For every fixed $a\in[0,1]$, definition \eqref{eq:Gamma} gives
\[
    \sum_{t=1}^T r_t^p(a)
    \le
    \sum_{t=1}^T r_{t,H}^p(a)
    +\GammaTail_{T,H}(p).
\]
Taking the supremum over $a$ yields
\[
    \sup_a\sum_{t=1}^T r_t^p(a)
    \le
    \sup_a\sum_{t=1}^T r_{t,H}^p(a)
    +\GammaTail_{T,H}(p).
\]
Moreover, $r_t^p(a)\ge r_{t,H}^p(a)$ for every $t,a$, so pathwise
\[
    \sum_{t=1}^T r_t^p(A_t)
    \ge
    \sum_{t=1}^T r_{t,H}^p(A_t).
\]
Taking expectations and combining the two inequalities proves
\eqref{eq:truncation}.
\end{proof}

\subsection{Discretization}

For $K\ge2$, let
\begin{equation}\label{eq:grid}
    \G_K
    :=
    \left\{0,\frac1K,\ldots,\frac{K-1}{K}\right\}.
\end{equation}

\begin{lemma}[Downward discretization]\label{lem:discretization}
Fix $T,H,p$, and suppose that the learner always chooses $A_t\in\G_K$.  Define
its grid-comparator truncated regret by
\[
    R_{T,H,K}(\pi,p)
    :=
    \max_{b\in\G_K}
    \sum_{t=1}^T r_{t,H}^p(b)
    -
    \E_\pi\!\left[\sum_{t=1}^T r_{t,H}^p(A_t)\right].
\]
Then
\begin{equation}\label{eq:disc}
    R_{T,H}(\pi,p)
    \le
    R_{T,H,K}(\pi,p)+\frac{T}{K}.
\end{equation}
\end{lemma}

\begin{proof}
For $a\in[0,1]$, define its downward grid approximation
\[
    b(a)
    :=
    \frac{\min\{\lfloor Ka\rfloor,K-1\}}{K}.
\]
Then $b(a)\in\G_K$, $b(a)\le a$, and
\[
    0\le a-b(a)\le\frac1K.
\]
Fix $a$ and write $b=b(a)$.  Since $b\le a$, every future price that exceeds
$a$ also exceeds $b$.  Hence
\[
    \ind{\Tau_t^p(a)\le H}
    \le
    \ind{\Tau_t^p(b)\le H}.
\]
Therefore, pathwise,
\begin{align*}
    r_{t,H}^p(a)
    &=a\ind{\Tau_t^p(a)\le H}\\
    &=b\ind{\Tau_t^p(a)\le H}
      +(a-b)\ind{\Tau_t^p(a)\le H}\\
    &\le b\ind{\Tau_t^p(b)\le H}+\frac1K\\
    &=r_{t,H}^p(b)+\frac1K.
\end{align*}
Summing over $t$ and taking the supremum over $a$ gives
\[
    \sup_{a\in[0,1]}\sum_{t=1}^T r_{t,H}^p(a)
    \le
    \max_{b\in\G_K}\sum_{t=1}^T r_{t,H}^p(b)
    +\frac{T}{K}.
\]
Subtracting the learner's expected truncated reward proves \eqref{eq:disc}.
\end{proof}

\subsection{A self-contained Hedge bound}

We first record Hoeffding's lemma in the form needed below.

\begin{lemma}[Hoeffding's lemma]\label{lem:hoeffding}
If $X\in[0,1]$ almost surely, then for every $\eta\in\R$,
\[
    \log\E[e^{\eta X}]
    \le
    \eta\E[X]+\frac{\eta^2}{8}.
\]
\end{lemma}

\begin{proof}
Let
\[
    \phi(\eta)
    :=
    \log\E[e^{\eta X}]-\eta\E[X].
\]
Under the exponentially tilted distribution proportional to $e^{\eta X}$,
$\phi''(\eta)$ is the variance of $X$.  Every random variable supported on
$[0,1]$ has variance at most $1/4$, so $\phi''(\eta)\le1/4$.  Also
$\phi(0)=\phi'(0)=0$.  Integrating the second-derivative bound twice gives
$\phi(\eta)\le\eta^2/8$.
\end{proof}

\begin{lemma}[Hedge with full-information gains]\label{lem:hedge}
Let $x_1,\ldots,x_n\in[0,1]^K$ be arbitrary gain vectors.  There is a
randomized online algorithm which chooses $I_s\in\{1,\ldots,K\}$ before seeing
$x_s$, observes all of $x_s$ afterward, and satisfies
\begin{equation}\label{eq:hedge}
    \max_{i\in\{1,\ldots,K\}}\sum_{s=1}^n x_{s,i}
    -
    \E\!\left[\sum_{s=1}^n x_{s,I_s}\right]
    \le
    \sqrt{\frac{n\log K}{2}}.
\end{equation}
For $n=0$, the regret is zero.
\end{lemma}

\begin{proof}
Assume $n\ge1$.  Initialize $w_{1,i}=1$ for every $i$.  At local round~$s$,
play arm $i$ with probability
\[
    q_{s,i}=\frac{w_{s,i}}{\sum_jw_{s,j}},
\]
and, after observing $x_s$, update
\[
    w_{s+1,i}=w_{s,i}e^{\eta x_{s,i}}.
\]
Let $W_s=\sum_iw_{s,i}$.  By Lemma~\ref{lem:hoeffding}, applied to a random
variable that equals $x_{s,i}$ with probability $q_{s,i}$,
\[
    \log\frac{W_{s+1}}{W_s}
    =
    \log\sum_iq_{s,i}e^{\eta x_{s,i}}
    \le
    \eta\sum_iq_{s,i}x_{s,i}+\frac{\eta^2}{8}.
\]
Summing over $s$ gives
\[
    \log\frac{W_{n+1}}{W_1}
    \le
    \eta\sum_{s=1}^n\sum_iq_{s,i}x_{s,i}
    +\frac{n\eta^2}{8}.
\]
For every fixed arm~$i$,
\[
    W_{n+1}
    \ge
    w_{n+1,i}
    =
    \exp\!\left(\eta\sum_{s=1}^n x_{s,i}\right),
\]
while $W_1=K$.  Hence
\[
    \sum_{s=1}^n x_{s,i}
    -
    \sum_{s=1}^n\sum_jq_{s,j}x_{s,j}
    \le
    \frac{\log K}{\eta}+\frac{n\eta}{8}.
\]
The expected gain of the randomized algorithm at round~$s$ is
$\sum_jq_{s,j}x_{s,j}$.  Maximizing over~$i$ and choosing
\[
    \eta=\sqrt{\frac{8\log K}{n}}
\]
proves \eqref{eq:hedge}.
\end{proof}

\subsection{Learning the truncated problem}

\begin{lemma}[Interlaced Hedge]\label{lem:interlaced}
For every $T,H,K$, there exists a nonanticipating grid-valued policy $\pi_{H,K}$
such that, for every deterministic sequence $p\in[0,1]^T$,
\begin{equation}\label{eq:interlaced}
    R_{T,H,K}(\pi_{H,K},p)
    \le
    \sqrt{\frac{T(H+1)\log K}{2}}.
\end{equation}
\end{lemma}

\begin{proof}
Partition the rounds into $H+1$ residue classes modulo $H+1$:
\[
    S_j
    :=
    \{t\in\{1,\ldots,T\}:t\equiv j\pmod{H+1}\},
    \qquad j=1,\ldots,H+1,
\]
using any consistent convention for residues.  Run one independent Hedge
instance on each $S_j$, with arms indexed by the quotes in $\G_K$.

For a round $t$, define the full-information gain vector
\[
    x_t(b):=r_{t,H}^p(b),\qquad b\in\G_K.
\]
This vector is completely determined after observing
\[
    p_{t+1},\ldots,p_{\min\{t+H,T\}}.
\]
If $t'$ is the next round after $t$ in the same residue class, then
$t'=t+H+1$.  Thus $x_t$ is known no later than round $t'-1$, before the next
decision of the same Hedge instance.  Consequently, each residue-class
instance receives ordinary immediate full-information feedback on its own local
time scale.

Let $n_j=|S_j|$.  Lemma~\ref{lem:hedge} gives local expected regret at most
$\sqrt{n_j\log K/2}$.  Since
\[
    \max_{b\in\G_K}\sum_{t=1}^T x_t(b)
    \le
    \sum_{j=1}^{H+1}
    \max_{b\in\G_K}\sum_{t\in S_j}x_t(b),
\]
the global grid-comparator regret is at most the sum of the local regrets.
By Cauchy--Schwarz,
\begin{align*}
    R_{T,H,K}(\pi_{H,K},p)
    &\le
    \sum_{j=1}^{H+1}\sqrt{\frac{n_j\log K}{2}}\\
    &\le
    \sqrt{(H+1)\left(\sum_{j=1}^{H+1}n_j\right)\frac{\log K}{2}}\\
    &=
    \sqrt{\frac{T(H+1)\log K}{2}}.
\end{align*}
\end{proof}

\subsection{Proof of the upper bound}

\begin{proof}[Proof of Theorem~\ref{thm:main}(i)]
Fix $H\ge0$ and $K\ge2$, and use the policy $\pi_{H,K}$ from
Lemma~\ref{lem:interlaced}.  For every $p\in\C_T(g)$, Lemmas
\ref{lem:truncation}, \ref{lem:discretization}, and \ref{lem:interlaced} give
\begin{align*}
    R_T(\pi_{H,K},p)
    &\le
    R_{T,H}(\pi_{H,K},p)+\GammaTail_{T,H}(p)\\
    &\le
    R_{T,H,K}(\pi_{H,K},p)
    +\frac{T}{K}
    +\GammaTail_{T,H}(p)\\
    &\le
    \sqrt{\frac{T(H+1)\log K}{2}}
    +\frac{T}{K}
    +Tg(H).
\end{align*}
The right-hand side is uniform over $p\in\C_T(g)$.  Taking the supremum over
$p$ and then the infimum over all policies proves \eqref{eq:upper-main}.
\end{proof}

\section{Lower bound}

The lower bound uses sequences whose only positive price occurs at the terminal
round.

\begin{lemma}[Tail mass of a terminal jump]\label{lem:terminal-tail}
For $x\in[0,1]$, let
\[
    p^x=(0,\ldots,0,x)\in[0,1]^T.
\]
Then, for every integer $H\ge0$,
\begin{equation}\label{eq:terminal-gamma}
    \GammaTail_{T,H}(p^x)
    =
    x\,(T-H-1)_+,
\end{equation}
where $(z)_+=\max\{z,0\}$.
\end{lemma}

\begin{proof}
Fix a quote $a$.  If $a\ge x$, it never executes.  If $a<x$, then for every
$t<T$ its first execution occurs at time~$T$, so
\[
    \Tau_t^{p^x}(a)=T-t.
\]
The delay is strictly larger than $H$ exactly for
$t=1,\ldots,T-H-1$, when this set is nonempty.  Hence the late-trade reward of
$a<x$ is
\[
    a\,(T-H-1)_+.
\]
Taking the supremum over $a<x$ gives $x(T-H-1)_+$.  The value need not be
attained because a quote equal to $x$ does not execute.
\end{proof}

\begin{lemma}[Membership of terminal jumps]\label{lem:terminal-membership}
Let $c=\lim_{H\to\infty}g(H)>0$.  Then
\[
    p^x\in\C_T(g)
    \qquad\text{for every }x\in[0,c].
\]
\end{lemma}

\begin{proof}
Because $g$ is nonincreasing with limit $c$, we have $g(H)\ge c$ for every
$H$.  By Lemma~\ref{lem:terminal-tail}, for $x\le c$,
\[
    \GammaTail_{T,H}(p^x)
    =x(T-H-1)_+
    \le xT
    \le cT
    \le Tg(H).
\]
This is precisely the defining condition of $\C_T(g)$.
\end{proof}

\begin{lemma}[Hard distribution over terminal prices]\label{lem:hard-prior}
Fix $c\in(0,1]$.  There exists a random variable $X\in[0,c]$ such that
\begin{equation}\label{eq:hard-properties}
    \E[X]=\frac{2c}{e}
    \qquad\text{and}\qquad
    a\,\Pr(X>a)\le\frac{c}{e}
    \quad\text{for every }a\in[0,1].
\end{equation}
\end{lemma}

\begin{proof}
Define $X$ through its survival function
\[
    \Pr(X>x)
    =
    \begin{cases}
        1, & 0\le x<c/e,\\[1mm]
        \dfrac{c}{ex}, & c/e\le x<c,\\[2mm]
        0, & x\ge c.
    \end{cases}
\]
Equivalently, $X$ has density $c/(ex^2)$ on $[c/e,c)$ and an atom of mass
$1/e$ at $c$.  Using the tail-integral formula,
\begin{align*}
    \E[X]
    &=\int_0^c\Pr(X>x)\,dx\\
    &=\frac{c}{e}
      +\int_{c/e}^c\frac{c}{ex}\,dx\\
    &=\frac{c}{e}+\frac{c}{e}
    =\frac{2c}{e}.
\end{align*}
If $a<c/e$, then $a\Pr(X>a)=a<c/e$.  If $c/e\le a<c$, then
$a\Pr(X>a)=c/e$.  If $a\ge c$, the product is zero.  This proves
\eqref{eq:hard-properties}.
\end{proof}

\begin{proof}[Proof of Theorem~\ref{thm:main}(ii)]
Fix an arbitrary randomized policy $\pi$.  Draw $X$ according to
Lemma~\ref{lem:hard-prior}, and present the terminal-jump sequence $p^X$.
By Lemma~\ref{lem:terminal-membership}, every realization belongs to
$\C_T(g)$.

For each $t<T$, the observed history under every realization of $X$ is the
all-zero prefix.  Therefore $A_t$ is independent of~$X$.  The learner's reward
from round~$t$ is
\[
    A_t\ind{A_t<X}.
\]
Conditioning on $A_t$ and applying \eqref{eq:hard-properties},
\[
    \E\bigl[A_t\ind{A_t<X}\bigr]
    =
    \E\bigl[A_t\Pr(X>A_t\mid A_t)\bigr]
    \le
    \frac{c}{e}.
\]
There is no reward from a quote submitted at time~$T$.  Hence
\begin{equation}\label{eq:learner-hard}
    \E\!\left[\sum_{t=1}^T r_t^{p^X}(A_t)\right]
    \le
    \frac{c}{e}(T-1).
\end{equation}

For a realized terminal value $X=x$, every quote $a<x$ posted before time~$T$
executes at time~$T$.  Thus the best fixed comparator value is
\[
    \sup_{a<x}(T-1)a=(T-1)x.
\]
Averaging over $X$ and using \eqref{eq:hard-properties},
\begin{equation}\label{eq:comparator-hard}
    \E\!\left[
        \sup_{a\in[0,1]}
        \sum_{t=1}^T r_t^{p^X}(a)
    \right]
    =(T-1)\E[X]
    =\frac{2c}{e}(T-1).
\end{equation}
Subtracting \eqref{eq:learner-hard} from \eqref{eq:comparator-hard} shows that
$\pi$ has average regret at least $c(T-1)/e$ under this distribution over
deterministic sequences.  Therefore at least one realization $p^x$ satisfies
\[
    R_T(\pi,p^x)\ge\frac{c}{e}(T-1).
\]
Since this holds for every policy $\pi$ and every such $p^x$ lies in
$\C_T(g)$, taking the infimum over policies proves \eqref{eq:lower-main}.
\end{proof}

\section{Deriving the if-and-only-if statement}

\begin{proof}[Proof of \eqref{eq:iff}]
Suppose first that $g(H)\to0$.  For $T\ge2$, choose
\[
    H_T=\lfloor T^{1/3}\rfloor,
    \qquad
    K_T=T.
\]
The upper bound \eqref{eq:upper-main} gives
\[
    \frac{\Reg_T(g)}{T}
    \le
    \sqrt{\frac{(H_T+1)\log T}{2T}}
    +\frac1T
    +g(H_T).
\]
Here $H_T\to\infty$ and $H_T\log T/T\to0$, so all three terms converge to
zero.  Therefore $\Reg_T(g)=o(T)$.

Conversely, if $g(H)$ does not converge to zero, then monotonicity implies that
$c=\lim_Hg(H)>0$.  By \eqref{eq:lower-main},
\[
    \liminf_{T\to\infty}\frac{\Reg_T(g)}{T}
    \ge
    \frac{c}{e}>0.
\]
Thus the minimax regret is not sublinear.
\end{proof}

\section{Remarks on the scope of the characterization}

\begin{remark}[Why maximality matters]
The iff statement is for the maximal class $\C_T(g)$ containing every sequence
that satisfies the tail envelope.  A smaller structured class can have very long
delays and still be easy.  For example, if the learner is told in advance that
the unique possible sequence is $(0,\ldots,0,1)$, it can post quotes arbitrarily
close to~$1$ from the first round.  Thus a delay-tail condition alone cannot
characterize every arbitrary subclass of deterministic sequences; it becomes
necessary precisely because the maximal class includes all hard terminal-jump
continuations allowed by the envelope.
\end{remark}

\begin{remark}[Interpretation of $\GammaTail_{T,H}$]
For the terminal-jump sequence $p^x$, Lemma~\ref{lem:terminal-tail} gives
\[
    \GammaTail_{T,H}(p^x)=x(T-H-1)_+.
\]
The final $H+1$ submission times do not count because their waiting times are at
most~$H$ or there is no future round.  The expression is nevertheless of order
$xT$ whenever $H=o(T)$.
\end{remark}

\end{document}
